\begin{table}[htbp]
\centering
\caption{Adversarial-Adaptation Simulation: AUC Degradation Under Strategic Cartel Adaptations}
\label{tab:adversarial_adaptation}
\footnotesize
\begin{adjustbox}{max width=\textwidth,center}
\begin{tabular}{p{8cm}ccc}
\toprule
Scenario & AUC & $\Delta$ vs baseline & Retention \\
\midrule
Baseline (no adaptation) & 0.722 & +0.000 & 100.0\% \\
Rotation: 20% of cover bidders retire and are replaced & 0.730 & +0.008 & 101.1\% \\
Occasional wins: 5% of FL14 firms get one win & 0.720 & -0.001 & 99.8\% \\
CNPJ splitting: each FL14 firm splits into two CNPJs & 0.653 & -0.068 & 90.6\% \\
Threshold cap: firms cap participation at 13 tender-items & 0.676 & -0.046 & 93.7\% \\
Combined: threshold cap + 20% rotation & 0.607 & -0.114 & 84.1\% \\
\bottomrule
\end{tabular}
\end{adjustbox}

\smallskip
{\footnotesize\textit{Notes:} Each scenario applies a strategic
adaptation to the BEC always-loser pool, holding the construct
($\log(1{+}\text{tenders\_count})$ ranking against $193$ CADE
cobidder labels) fixed. The construct is resilient to rotation
and to occasional wins (the bright-line filter eliminates the
small share of newly-winning cover bidders without compressing
the ranking), and degrades modestly under CNPJ splitting. It is
substantially vulnerable to threshold-aware capping and severely
degraded when capping is combined with rotation: a sophisticated
cartel that targets the threshold can compress AUC by roughly
$30$ points. Retention is the fraction of baseline AUC preserved.\par}
\end{table}
